\documentclass[14pt]{extarticle}
\usepackage{fontspec}
\usepackage{polyglossia}
\usepackage{amsmath}
\usepackage{amssymb}
\usepackage{xcolor}
\usepackage{fancyhdr}
\usepackage{graphicx}
\usepackage{listings}
\usepackage[most]{tcolorbox}
% \usepackage{setspace}
\usepackage{pifont}
\usepackage{enumitem}
\usepackage{titlesec}
\usepackage[bottom]{footmisc}
\usepackage{titling}
\usepackage{minted}
\usepackage{etoolbox}
\usepackage{array}
\usepackage{longtable} % page-breakable tables (must load before polyglossia/bidi RTL setup)
\usepackage{extsizes}
\usepackage{float}
\usepackage{booktabs}
\usepackage{hyperref}
\usepackage[onehalfspacing]{setspace}
\newif\ifwhiteblack
\ifdefined\setwhiteblack
  \setwhiteblack
\fi

\ifwhiteblack
  \usepackage[a4paper, landscape, top=1.25cm, bottom=1.25cm, left=1.25cm, right=1.25cm, headheight=1.5cm, headsep=0.5cm]{geometry}
\else
  \usepackage[a4paper, top=2.7cm, bottom=2.7cm, left=2.5cm, right=2.5cm, headheight=1.5cm, headsep=0.5cm]{geometry}
\fi
\usepackage{tikz}
\usepackage{datetime2}

\usetikzlibrary{automata, positioning, arrows.meta, arrows}
\providecommand{\LTcaptype}{table}

\newfontfamily\emoji{Segoe UI Emoji}

\pagestyle{fancy}

\setmainlanguage[numerals=western]{arabic}
\setotherlanguage{english}
\newfontfamily\arabicfont[Script=Arabic]{Calibri}
\newfontfamily\arabicfonttt[Script=Arabic]{Courier New}



\ifwhiteblack
\lstset{
  language=[Sharp]C,
  numbers=left,
  stepnumber=1,
  numbersep=8pt,
  frame=single,
  basicstyle=\ttfamily\small,
  keywordstyle=\bfseries,
  stringstyle=\ttfamily,
  commentstyle=\itshape
}
\else
\lstset{
  language=[Sharp]C,
  numbers=left,
  stepnumber=1,
  numbersep=8pt,
  frame=single,
  basicstyle=\ttfamily\small,
  keywordstyle=\color{blue},
  stringstyle=\color{red},
  commentstyle=\color{green!50!black}
}
\fi

\newif\ifdetailed
\ifdefined\setdetailed
  \setdetailed
\fi

\newif\ifwithsols
\ifdefined\setwithsols
  \setwithsols
\fi

\ifwhiteblack
  \AtBeginDocument{
    \LARGE
    \usemintedstyle{bw}
  }
\fi

% Default minted look: compact code that wraps long lines.
% Not applied in black-and-white landscape mode, which already uses \LARGE.
\ifwhiteblack\else
  \setminted{fontsize=\small, breaklines=true}
\fi

% unified tcolorboxes for articles
\ifwhiteblack
\tcbset{colback=white, colframe=black, fonttitle=\bfseries, boxrule=1pt}
\newtcolorbox{boxDef}[1][]{title={{\emoji📘} تعريف\ifx\\#1\\\else ~#1\fi :}}
\newtcolorbox{boxExercise}[1][]{title={{\emoji🧩} تمرين\ifx\\#1\\\else ~#1\fi :}}
\newtcolorbox{boxExample}[1][]{title={{\emoji📝} مثال\ifx\\#1\\\else ~#1\fi :}}
\newtcolorbox{boxNote}[1][]{title={{\emoji✨} ملاحظة\ifx\\#1\\\else ~#1\fi :}}
\newtcolorbox{boxAttention}[1][]{title={{\emoji🔔} تنبيه\ifx\\#1\\\else ~#1\fi :}}
\newtcolorbox{boxWarning}[1][]{title={{\emoji⚡} ملاحظة هامة\ifx\\#1\\\else ~#1\fi :}}
\newtcolorbox{boxSolution}[1][]{title={{\emoji✅} حل\ifx\\#1\\\else ~#1\fi :}}
\newtcolorbox{boxSymbol}[1][]{title={{\emoji🔣} رمز\ifx\\#1\\\else ~#1\fi :}}
\newtcolorbox{boxHint}[1][]{title={{\emoji💡} تلميح\ifx\\#1\\\else ~#1\fi :}}
\else
\tcbset{colback=white, colframe=black, fonttitle=\bfseries, boxrule=0.8pt}
\newtcolorbox{boxDef}[1][]{colback=blue!5!white,colframe=blue!75!black,
  title={{\emoji📘} تعريف\ifx\\#1\\\else ~#1\fi :}}
\newtcolorbox{boxExercise}[1][]{colback=cyan!5!white,colframe=cyan!70!black,
  title={{\emoji🧩} تمرين\ifx\\#1\\\else ~#1\fi :}}
\newtcolorbox{boxExample}[1][]{colback=yellow!5!white,colframe=orange!90!black,
  title={{\emoji📝} مثال\ifx\\#1\\\else ~#1\fi :}}
\newtcolorbox{boxNote}[1][]{colback=gray!10!white,colframe=black,
  title={{\emoji✨} ملاحظة\ifx\\#1\\\else ~#1\fi :}}
\newtcolorbox{boxAttention}[1][]{colback=magenta!10!white,colframe=magenta!80!black,
  title={{\emoji🔔} تنبيه\ifx\\#1\\\else ~#1\fi :}}
\newtcolorbox{boxWarning}[1][]{colback=red!5!white,colframe=red!75!black,
  title={{\emoji⚡} ملاحظة هامة\ifx\\#1\\\else ~#1\fi :}}
\newtcolorbox{boxSolution}[1][]{colback=green!5!white,colframe=green!60!black,
  title={{\emoji✅} حل\ifx\\#1\\\else ~#1\fi :}}
\newtcolorbox{boxSymbol}[1][]{colback=purple!5!white,colframe=purple!70!black,
  title={{\emoji🔣} رمز\ifx\\#1\\\else ~#1\fi :}}
\newtcolorbox{boxHint}[1][]{colback=teal!5!white,colframe=teal!60!black,
  title={{\emoji💡} تلميح\ifx\\#1\\\else ~#1\fi :}}
\fi


\ifwhiteblack
\tcbset{simplecode/.style={ colback=white, colframe=black, boxrule=1pt, arc=2pt, left=4pt,right=4pt,top=4pt,bottom=4pt}}
\else
\tcbset{simplecode/.style={ colback=gray!5, colframe=black!50, boxrule=0.4pt, arc=2pt, left=4pt,right=4pt,top=4pt,bottom=4pt}}
\fi
% Compact code boxes: tighter line spacing (the preamble uses onehalfspacing) and small vertical gaps.
% Must come AFTER the simplecode definition above.
\tcbset{simplecode/.append style={fontupper=\setstretch{0.95}, top=4pt, bottom=4pt, before skip=6pt, after skip=6pt}}
\newenvironment{boxCode}{\begin{tcolorbox}[simplecode]}{\end{tcolorbox}}

% Code cell used inside tables (keeps code left-to-right inside RTL tables).
% Use inside a fixed-width p{...} column: \begin{pycell}\begin{minted}...\end{minted}\end{pycell}
\newtcolorbox{pycell}{frame empty,colback=white,boxsep=0pt,left=2pt,right=2pt,top=2pt,bottom=2pt,width=\linewidth}

\newcolumntype{C}[1]{>{\centering\arraybackslash}p{#1}}

\makeatletter
\DTMnewdatestyle{dotted}{%
  \renewcommand{\DTMdisplaydate}[4]{\two@digits{##3}.\two@digits{##2}.##1}%
  \renewcommand{\DTMDisplaydate}{\DTMdisplaydate}%
}
\makeatother
\DTMsetdatestyle{dotted}

% redefine spaces after titles
\makeatletter
\renewcommand{\@maketitle}{%
  \begin{center}
    {\huge \bfseries \@title \par}%
    \vskip 0.2em % space between title and author
    {\large \@author \par}%
    % \vskip 0.2em % space between author and date
    % {\normalsize \@date \par}%
    {\small \hfill \begin{english}\DTMtoday\end{english} \par}%
  \end{center}
}
\makeatother

\fancyhf{} % clear default
\fancypagestyle{plain}{
  \fancyhf{}
  \fancyhead[L]{مدرسة التسامح الشاملة}
  % \fancyhead[L]{\includegraphics[height=1cm]{../../../images/logoTasamoh.png}}
  \fancyhead[R]{الأستاذ محمود اغبارية}
  \fancyfoot[C]{\thepage}
}

\fancyhead[L]{مدرسة التسامح الشاملة}
\fancyhead[R]{الأستاذ محمود اغبارية}
\fancyfoot[C]{\thepage}

\setcounter{tocdepth}{3} % only section subsection and subsubsection in TOC

\newcommand{\importpdfpage}[2]{%
    \thispagestyle{empty}
    \newgeometry{left=0cm,right=0cm,top=0cm,bottom=0cm}
    \includegraphics[page=#2,width=0.95\paperwidth,height=0.95\paperheight,keepaspectratio]{#1}
    \restoregeometry
}

\newcommand{\insertFullImg}[1]{%
    \noindent
    \makebox[\textwidth][c]{
        \includegraphics[width=0.9\paperwidth,keepaspectratio]{#1}%
    }%
}%


\ifwithsols
\title{حل ورقة عمل - تدريب على الشروط والحلقات (بايثون أ)}
\else
\title{ورقة عمل - تدريب على الشروط والحلقات (بايثون أ)}
\fi

\begin{document}

\maketitle
\thispagestyle{fancy}

\section*{تمارين: الشروط والحلقات}

\begin{enumerate}[label=\textbf{تمرين \arabic*:}, wide=0pt, itemsep=2.2em]

% ---------------------------------------------------------------
\item طلبت المعلمة من الطلاب تنفيذ المهمة التالية:
"اكتبوا برنامجًا يستقبل عددين صحيحين، ويطبع عبارة توضيحية تصف إذا كان العددان \textbf{زوجيين} أو \textbf{فرديين}".

كتب الطلاب عدة حلول مختلفة. افترضوا أنّه في بداية كل مقطع برنامج كتب الطلاب السطرين التاليين:
\begin{boxCode}
\begin{english}
\begin{minted}{python}
num1 = int(input('Enter an integer number: '))
num2 = int(input('Enter an integer number: '))
\end{minted}
\end{english}
\end{boxCode}

حددوا لكل مقطع برنامج إذا كان يحقّق المطلوب (نعم أم لا)؛ وإذا كان مقطع البرنامج لا يحقّق المطلوب، اشرحوا لماذا.

\begin{enumerate}[label=\alph*., itemsep=1.5em]

\item .
\begin{boxCode}
\begin{english}
\begin{minted}{python}
if num1 % num2 == 0:
    print(num1, num2, "Both are even")
else:
    print(num1, num2, "Both are odd")
\end{minted}
\end{english}
\end{boxCode}
\ifwithsols
\begin{boxSolution}
\textbf{لا يحقّق المطلوب.} لأنّ الشرط \textenglish{\texttt{num1 \% num2 == 0}} يفحص إذا كان \textenglish{num1} يقسم على \textenglish{num2} بدون باقٍ، ولا يفحص إذا كان كل عدد يقسم على $2$ بدون باقٍ. مثلًا لـ $\text{num1} = 9$ و $\text{num2} = 3$ يطبع أنّهما زوجيان رغم أنّهما فرديان.
\end{boxSolution}
\else
\textbf{هل يحقّق المطلوب (نعم / لا)؟} \underline{\hspace{3cm}} \quad \textbf{الشرح:} \underline{\hspace{7cm}}
\fi
\filbreak

\item .
\begin{boxCode}
\begin{english}
\begin{minted}{python}
even = num1
odd = num2
print(num1, "Even number")
print(num2, "Odd number")
\end{minted}
\end{english}
\end{boxCode}
\ifwithsols
\begin{boxSolution}
\textbf{لا يحقّق المطلوب.} لأنّ مقطع البرنامج لا يفحص قيمتي العددين أصلًا، بل يفترض دائمًا أنّ العدد الأول زوجي والعدد الثاني فردي بغضّ النظر عن المُدخل.
\end{boxSolution}
\else
\textbf{هل يحقّق المطلوب (نعم / لا)؟} \underline{\hspace{3cm}} \quad \textbf{الشرح:} \underline{\hspace{7cm}}
\fi
\filbreak

\item .
\begin{boxCode}
\begin{english}
\begin{minted}{python}
if num1 % 2 == 0:
    print(num1, num2, "Both are even")
else:
    print(num1, num2, "Both are odd")
\end{minted}
\end{english}
\end{boxCode}
\ifwithsols
\begin{boxSolution}
\textbf{لا يحقّق المطلوب.} لأنّ الشرط يفحص العدد الأول \textenglish{num1} فقط ويتجاهل فحص العدد الثاني \textenglish{num2}. مثلًا لـ $\text{num1} = 4$ و $\text{num2} = 5$ يطبع أنّهما زوجيان.
\end{boxSolution}
\else
\textbf{هل يحقّق المطلوب (نعم / لا)؟} \underline{\hspace{3cm}} \quad \textbf{الشرح:} \underline{\hspace{7cm}}
\fi
\filbreak

\item .
\begin{boxCode}
\begin{english}
\begin{minted}{python}
if num1 % 2 == 0 and num2 % 2 == 0:
    print(num1, num2, "Both are even")
else:
    print(num1, num2, "Both are odd")
\end{minted}
\end{english}
\end{boxCode}
\ifwithsols
\begin{boxSolution}
\textbf{لا يحقّق المطلوب.} لأنّ فرع \textenglish{else} يعمل عندما لا يكون العددان زوجيين معًا، وهذا يشمل الحالة التي يكون فيها أحدهما زوجيًا والآخر فرديًا (مثل $\text{num1} = 4$ و $\text{num2} = 5$)، فيطبع بالخطأ أنّهما فرديان.
\end{boxSolution}
\else
\textbf{هل يحقّق المطلوب (نعم / لا)؟} \underline{\hspace{3cm}} \quad \textbf{الشرح:} \underline{\hspace{7cm}}
\fi
\filbreak

\item .
\begin{boxCode}
\begin{english}
\begin{minted}{python}
if num1 % 2 == 0 and num2 % 2 == 0:
    print(num1, num2, "Both are even")
elif num1 % 2 != 0 and num2 % 2 != 0:
    print(num1, num2, "Both are odd")
else:
    print("One is even and the other is odd")
\end{minted}
\end{english}
\end{boxCode}
\ifwithsols
\begin{boxSolution}
\textbf{نعم، يحقّق المطلوب.} لأنّه يفحص أولًا إذا كان العددان زوجيين معًا، ثم يفحص إذا كانا فرديين معًا، وفي باقي الحالات يطبع أنّ أحدهما زوجي والآخر فردي.
\end{boxSolution}
\else
\textbf{هل يحقّق المطلوب (نعم / لا)؟} \underline{\hspace{3cm}} \quad \textbf{الشرح:} \underline{\hspace{7cm}}
\fi
\filbreak

\item .
\begin{boxCode}
\begin{english}
\begin{minted}{python}
if num1 % 2 == 0:
    if num2 % 2 == 0:
        print(num1, num2, "Both are even")
    else:
        print("One is even and the other is odd")
elif num1 % 2 != 0 and num2 % 2 != 0:
    print(num1, num2, "Both are odd")
else:
    print("One is even and the other is odd")
\end{minted}
\end{english}
\end{boxCode}
\ifwithsols
\begin{boxSolution}
\textbf{نعم، يحقّق المطلوب.} لأنّه يغطي جميع الحالات الأربعة الممكنة لزوجية وفردية العددين بشكل صحيح بواسطة شروط متداخلة ومتتالية.
\end{boxSolution}
\else
\textbf{هل يحقّق المطلوب (نعم / لا)؟} \underline{\hspace{3cm}} \quad \textbf{الشرح:} \underline{\hspace{7cm}}
\fi
\filbreak

\end{enumerate}
\filbreak

% ---------------------------------------------------------------
\item أمامكم جدول به مقطعا برنامج. لكل من المقطعين، اكتبوا ما هو المُخرج الناتج في العمود الفارغ.

{\renewcommand{\arraystretch}{1.4}
\begin{longtable}{|C{1.5cm}|p{8.5cm}|C{4.5cm}|}
\hline
\textbf{البند} & \centering\textbf{الكود بلغة بايثون} & \textbf{المُخرج الناتج} \\
\hline
أ &
\begin{pycell}
\begin{english}
\begin{minted}{python}
for i in range(-5, 0, 1):
    print(i)
\end{minted}
\end{english}
\end{pycell}
& \ifwithsols \begin{english}\texttt{-5}\newline\texttt{-4}\newline\texttt{-3}\newline\texttt{-2}\newline\texttt{-1}\end{english} \fi \\
\hline
ب &
\begin{pycell}
\begin{english}
\begin{minted}{python}
for i in range(5):
    print(i)
else:
    print("Done!")
\end{minted}
\end{english}
\end{pycell}
& \ifwithsols \begin{english}\texttt{0}\newline\texttt{1}\newline\texttt{2}\newline\texttt{3}\newline\texttt{4}\newline\texttt{Done!}\end{english} \fi \\
\hline
\end{longtable}}
\filbreak

% ---------------------------------------------------------------
\item أمامكم مقطع البرنامج التالي:

{\renewcommand{\arraystretch}{1.4}
\begin{longtable}{|p{9.5cm}|C{5cm}|}
\hline
\centering\textbf{الكود بلغة بايثون} & \textbf{المُخرج الناتج} \\
\hline
\begin{pycell}
\begin{english}
\begin{minted}{python}
num = 98713
reverse_number = 0
while num > 0:
    reminder = num % 10
    reverse_number = (reverse_number * 10) + reminder
    num = num // 10
print("Reverse number: ", reverse_number)
\end{minted}
\end{english}
\end{pycell}
& \ifwithsols \begin{english}\texttt{Reverse number:\ \ 31789}\end{english} \fi \\
\hline
\end{longtable}}

\vspace{1em}
\begin{enumerate}[label=\alph*.]
\item اكتبوا ما هو المُخرج الناتج في العمود الفارغ في الجدول أعلاه.
\ifwithsols
\begin{boxSolution}
المُخرج الناتج هو: \textenglish{\texttt{Reverse number:\ \ 31789}}

\begin{center}
\begin{english}
\renewcommand{\arraystretch}{1.3}
\begin{tabular}{|C{2cm}|C{2cm}|C{3cm}|C{5cm}|}
\hline
\textbf{num} & \textbf{reminder} & \textbf{reverse\_number} & \textarabic{\textbf{المخرجات}} \\
\hline
98713 & & 0 & \\ \hline
9871 & 3 & 3 & \\ \hline
987 & 1 & 31 & \\ \hline
98 & 7 & 317 & \\ \hline
9 & 8 & 3178 & \\ \hline
0 & 9 & 31789 & \texttt{Reverse number:\ \ 31789} \\ \hline
\end{tabular}
\end{english}
\end{center}
\end{boxSolution}
\fi
\filbreak

\item ما هو الهدف من مقطع البرنامج؟
\ifwithsols
\begin{boxSolution}
يقوم مقطع البرنامج بعكس ترتيب منازل العدد الصحيح الموجب \textenglish{num} وحفظ العدد المقلوب في المتغير \textenglish{reverse\_number} ثم طباعته.
\end{boxSolution}
\fi
\filbreak
\end{enumerate}
\filbreak

\end{enumerate}

\vspace{1cm}
\begin{flushleft}
أرجو لكم وقتًا ممتعًا.

الأستاذ محمود اغبارية.
\end{flushleft}

\end{document}
